\documentclass[12pt,a4paper]{ctexart}
\usepackage{geometry}\geometry{margin=2.2cm}
\usepackage{graphicx,float,fancyhdr,hyperref,longtable,booktabs,enumitem,xcolor,listings}
\lstset{basicstyle=\ttfamily\small,breaklines=true,frame=single}
\pagestyle{fancy}\fancyhf{}\fancyhead[L]{dgiot\_lite 离线开发环境用户手册}\fancyhead[R]{V1.0}\fancyfoot[C]{\thepage}
\title{\textbf{dgiot\_lite 离线开发环境\\用户手册}}
\author{约翰·刘}
\date{2026年7月13日}

\begin{document}\maketitle
\begin{abstract}
本手册介绍 dgiot\_lite IO 服务器离线开发环境的安装、配置和使用方法。
该环境在本地完整模拟了油田 131 生产服务器的 IO 采集架构，
支持 CommBridge、IoMonitor、IoCommit、IoProject、pSpace、OPC DA、Modbus TCP 七项服务，
可离线进行协议测试、采集器开发和全链路验证。
\end{abstract}

\section{环境概述}

\subsection{模拟架构}
\begin{verbatim}
  生产环境 (11.66.12.131)          本机模拟 (127.0.0.1)
  ──────────────────────          ──────────────────
  CommBridge  :53001       →      :53002  DTU桥接
  IoMonitor                →      :9002   数据汇聚
  IoCommit ×7              →      :9003   数据写入
  IoProject                →      :9001   进程编排
  pSpace :8889             →      :18889  pSpace数据源
  Kepware OPC :135         →      :13500  OPC DA数据源
  Modbus TCP :502          →      :502    直连轮询
\end{verbatim}

\subsection{服务端口清单}
\begin{table}[H]\centering
\begin{tabular}{lll}
\toprule
\textbf{服务} & \textbf{端口} & \textbf{协议} \\
\midrule
CommBridge & 53002 & DTU透传 (0xAA + ASCII\_ID) \\
IoMonitor & 9002 & JSON over TCP \\
IoCommit & 9003 & JSON写入 \\
IoProject & 9001 & JSON指令 (status/list) \\
IO-Srv (pSpace) & 18889 & pSpace模拟协议 \\
OPC-DA & 13500 & OPC模拟协议 \\
Modbus TCP & 502 & Modbus TCP (Func 0x03) \\
\bottomrule
\end{tabular}
\end{table}

\section{快速启动}

\subsection{一键启动全栈模拟}
\begin{lstlisting}
# 默认 20 台设备
python tools/dev_env.py

# 指定设备数量
python tools/dev_env.py --scale 56

# 只启动服务，不启动采集器
python tools/dev_env.py --no-collector
\end{lstlisting}

\textbf{启动后输出}：
\begin{verbatim}
  [CB] :53002
  [IoMon] :9002
  [IoCom] :9003
  [IoProj] :9001
  [IO-Srv] :18889
  [OPC] :13500
  [MB] :502

  验证连通性...
    CB:53002 OK    IoMon:9002 OK    IoCom:9003 OK
    IO-Srv:18889 OK  OPC:13500 OK    MB:502 OK

  RTU 在线: 20/20
\end{verbatim}

\subsection{交互控制台命令}
\begin{table}[H]\centering
\begin{tabular}{lp{5cm}}
\toprule
\textbf{命令} & \textbf{说明} \\
\midrule
\texttt{status} & 查看各服务运行状态 \\
\texttt{devices} & 列出在线设备及其轮询次数 \\
\texttt{perf} & 性能统计（报文速率、记录速率） \\
\texttt{inject <dev\_id> <type>} & 注入新设备（type: 00-B0） \\
\texttt{help} & 帮助 \\
\texttt{quit} & 退出 \\
\bottomrule
\end{tabular}
\end{table}

\section{pSpace 实时采集器}

\subsection{环境要求}
\begin{itemize}[leftmargin=*]
  \item 32位 Python 3.11（C:\textbackslash Python311-32\textbackslash python.exe）
  \item IO ServerOnLine 目录（含 psAPISDK.dll 及依赖）
  \item 网络可达 11.66.12.130:8889（仅真连需要）
\end{itemize}

\subsection{使用示例}
\begin{lstlisting}
# 读取指定 Tag ID
C:\Python311-32\python.exe tools/pspace_collector.py \
    --ids "5000,5501,6001,10000"

# 扫描活跃 Tag
C:\Python311-32\python.exe tools/pspace_collector.py \
    --scan --start 1 --end 20000 --step 100

# 默认读取（5000-5500范围）
C:\Python311-32\python.exe tools/pspace_collector.py
\end{lstlisting}

\subsection{输出示例}
\begin{verbatim}
Connected: handle=0x0001

Tag ID          Value      Timestamp              Quality
----------------------------------------------------------
5000           1.8750     2026-07-13T19:47:06    GOOD
5501           2.2437     2026-07-13T19:46:59    GOOD
6001           3.6946     2025-12-04T14:06:54    0x1C
10000          3.0440     2026-07-11T20:12:20    0x1C
\end{verbatim}

\section{协议测试工具}

\subsection{CommBridge 协议验证}
\begin{lstlisting}
# 本地 CommBridge 全链路测试
python tools/run_commbridge_local.py

# 输出: RTU注册→CommBridge轮询→IoMonitor→IoCommit
\end{lstlisting}

\subsection{Modbus TCP 直连测试}
\begin{lstlisting}
# 启动 Modbus mock 服务（:502）
python tools/mock_opc_server.py --port 502

# 客户端读取
python -c "
import socket, struct
s = socket.socket(); s.connect(('127.0.0.1',502))
s.send(struct.pack('>HHH',1,0,6)+b'\x01\x03\x00\x00\x00\x04')
resp = s.recv(256)
for i in range(4):
    v = struct.unpack('>H', resp[9+i*2:11+i*2])[0]
    print(f'Reg[{i}] = {v} -> {v*170/8192:.2f}A')
"
\end{lstlisting}

\subsection{OPC DA 模拟测试}
\begin{lstlisting}
python -c "
import socket, struct
s = socket.socket(); s.connect(('127.0.0.1',13500))
items = '02204060100.Ia;02204060100.Ua'
s.send(struct.pack('>HH',len(items)+4,0x0001)+items.encode())
resp = s.recv(512)
print(resp[4:].decode())
"
\end{lstlisting}

\section{Tag 扫描与点位发现}

\subsection{批量扫描 Tag ID}
\begin{lstlisting}
# 批量扫描 (子进程隔离, 自动限流)
python tools/scan_tags.py 60 1 20000 500

# Shell 快速边界探测
for base in 5000 5500 6000 ... 12000; do
  ids=$(python -c "print(','.join(str(base+i) for i in range(15)))")
  C:/Python311-32/python.exe tools/pspace_collector.py --ids "$ids" 2>/dev/null | head -3
done
\end{lstlisting}

\subsection{已知 Tag ID 图谱}
\begin{table}[H]\centering
\begin{tabular}{rlcl}
\toprule
\textbf{Block} & \textbf{典型值} & \textbf{时效} & \textbf{推断设备} \\
\midrule
5000 & 1.875 (平衡相) & 实时 & 02204060029 \\
5500 & 0--8.365 (混合) & 实时 & 02204060071 \\
7500 & 2.54--3.69 & 实时 & 02204060086 \\
8000 & 2.87--3.69 & 实时 & 02204060092 \\
9000 & 1.87--2.38 & 实时 & 02204060100 \\
9500 & 1.875 (平衡相) & 实时 & 02204060111 \\
10000 & 2.98--3.04 & OLD & 02204060117 \\
10500 & 2.94 & 实时 & 02204060121 \\
11000 & 1.99 & 实时 & 02204060123 \\
25000 & 2.75 (带负值) & 实时 & 待映射 \\
55000 & -6.18 (无功) & 实时 & 待映射 \\
\bottomrule
\end{tabular}
\end{table}

\textbf{规律}: 每设备占 $\sim$11-20 连续 Tag，块间距 500。值 1.875=平衡三相断路器，负值=无功功率 Q。

\subsection{限制与建议}
\begin{itemize}[leftmargin=*]
  \item 每次 Connect 仅能执行一次有效 ReadList
  \item 批量上限 $\sim$20 个 Tag
  \item 间隔需 30-60 秒避免限流
  \item 推荐改用 OPC Browse 直接浏览 Kepware tag tree
\end{itemize}

\section{设备配置}

\subsection{设备类型}
\begin{longtable}{lll}
\toprule
\textbf{编码} & \textbf{名称} & \textbf{通道数} \\
\midrule
0x00 & DSL-31A 断路器 & 20 \\
0x10 & DST-31A 变压器差动 & 15 \\
0x20 & DBPA-31A 备用电源自投 & 13 \\
0x30 & DSB-31A 变压器后备 & 20 \\
0x40 & 电动机保护 & 19 \\
0x50 & DST-22D 变压器差动 & 20 \\
0x60 & DSB-22D 变压器后备 & 20 \\
0x70 & DSL-24D 断路器 & 20 \\
0x80 & DGP-11 变压器差动保护 & 21 \\
0x90 & DGP-12 变压器后备保护 & 24 \\
0xA0 & DGP-13 接地保护 & 22 \\
0xB0 & DMP-31A 电动机保护 & 19 \\
\bottomrule
\end{longtable}

\subsection{转换系数}
\begin{table}[H]\centering
\begin{tabular}{lll}
\toprule
\textbf{系数} & \textbf{值} & \textbf{适用} \\
\midrule
C0 & 170/8192 & 电流/电压 (Ia,Ib,Ic,Ua,Ub,Uc) \\
C1 & 8.5/8192 & 接地电流 (Iac) \\
C2 & 170/8192 & 相电压 (Uphase) \\
C3 & 170×8.5/8192 & 有功功率 (P) \\
C4 & 1/8192 & 功率因数 (cos$\phi$) \\
C5 & 2/8192 & 频率 (F=50+Y×2/8192) \\
C6-9 & 1 & 状态量（直通） \\
\bottomrule
\end{tabular}
\end{table}

物理值 $Y = raw \times Coefficient[i]$，基准标定值 8192(0x2000)。

\section{数据校验规则}

\begin{table}[H]\centering
\begin{tabular}{lll}
\toprule
\textbf{层级} & \textbf{检查项} & \textbf{判定标准} \\
\midrule
L1 & 帧长度匹配 & abs(expected$-$actual) $\leq$ 2 \\
L2 & 值范围 & 电流 0-500A / 电压 100-400V / 功率 0-300kW \\
L3 & 三相平衡 & (I$_{max}$$-$I$_{min}$)/I$_{avg}$ $<$ 25\% \\
L4 & 历时一致性 & 相邻值变化 $<$ 50\% \\
L5 & Oracle 对标 & |A$-$B|/A $<$ 1\% \\
\bottomrule
\end{tabular}
\end{table}

\section{常见问题}

\begin{description}[leftmargin=*,style=nextline]
  \item[Q: psAPISDK 连接返回 -19998？]
  A: 表示"未初始化"。确保先调用 \texttt{psAPI\_Common\_StartAPI()}。
  \item[Q: Connect 返回 -18591？]
  A: 密码错误。确认使用正确的 pSpace 凭据。
  \item[Q: ReadList 返回 -19998？]
  A: 连接未正确建立，或 Tag ID 范围无效。重新 Connect 获取新 handle 后重试。
  \item[Q: 模拟环境端口被占用？]
  A: 使用 \texttt{netstat -ano | findstr ":端口"} 查找占用进程，\texttt{taskkill /PID xxx} 终止。
  \item[Q: 32位 Python 如何安装？]
  A: 下载 embeddable 版本，解压到 C:\textbackslash Python311-32，安装 pip 后即可使用。
\end{description}

\section{工具清单}

\begin{table}[H]\centering
\begin{tabular}{lp{8cm}}
\toprule
\textbf{工具} & \textbf{用途} \\
\midrule
\texttt{dev\_env.py} & 一键启动全栈模拟（7服务+N设备） \\
\texttt{pspace\_collector.py} & pSpace 实时数据采集器 \\
\texttt{mock\_opc\_server.py} & Modbus TCP 模拟服务器 \\
\texttt{simulate\_131\_fullscale.py} & 全量设备 1:1 模拟 \\
\texttt{run\_commbridge\_local.py} & CommBridge 全链路测试 \\
\texttt{probe\_pspace\_proto.py} & pSpace 协议探测+IO Server 模拟 \\
\texttt{cap\_opc\_local.py} & 本地 OPC 报文抓包代理 \\
\texttt{loop\_131\_monitor.py} & 131 生产服务器巡检 \\
\texttt{loop\_dgiot\_platform.py} & 本地平台健康巡检 \\
\texttt{platform\_auto\_restart.py} & 本地平台自动恢复 \\
\bottomrule
\end{tabular}
\end{table}

\end{document}
